\chapter*{Abstract}
\addcontentsline{toc}{chapter}{Abstract}
\markboth{Abstract}{}

\noindent Domestic intrusion alarms are now built almost entirely on single-chip Wi-Fi
systems-on-chip. On the ESP32, the radio driver, the TCP/IP stack, the MQTT client and the
alarm application share one die, one scheduler and one interrupt controller. Two pieces of
standard engineering advice are used to contain the resulting risk: pin the safety task to its
own core, and gate risky features behind build-time configuration. Both are assertions about
the system, and in practice neither is checked.

\vspace{6pt}
\noindent This work treats an ESP32 intrusion alarm as a case study and sets out to falsify
both assertions rather than to demonstrate them. Three defects were found, each residing in a
different kind of assumption, and each exposed by a method that costs considerably less than
the system it examines. A compile-time guard accepted the combination of TLS and a broker login
as evidence of a publisher's authority; under a single leaked credential the alarm fired in
only 16.7\% of 300 trials, against 100\% for the control. A response-time decomposition
exposed an unbounded command drain in the safety loop that 9,600 adversarial trials, 3,000,000
fuzzing iterations and exhaustive model checking had all passed over. The project's own
temporal-logic invariant forbidding unauthenticated remote disarm proved to be vacuous in the
specification that had been shipped and checked in continuous integration.

\vspace{6pt}
\noindent A twelve-configuration ablation over 3,600 trials shows that availability, integrity
and the safety function fail independently of one another: a token-bucket rate limiter costs
41\% of legitimate operator commands while contributing nothing to integrity. Broker access
control and application-layer message authentication were found to fail under different
compromises, a result reproduced against a live Mosquitto broker across eight scenarios. An
adversary in possession of a compromised controller host---holding both a credential the access
control list must permit and the command key---defeats the layered configuration and locks the
legitimate owner out of all 1,250 attempted commands.

\vspace{6pt}
\noindent An attempt to generalise the configuration finding across 135 commit-pinned
open-source embedded MQTT projects returned no hits for the guard pattern. That negative result
is reported as such, and no claim is made about the prevalence of the pattern in the wider
ecosystem.

\vspace{6pt}
\noindent The results establish \emph{logical} isolation only. Physical timing invariance is a
property of the silicon and of a closed-source radio driver; it has not been measured, no
latency figure is reported, and the hardware protocol that would settle the question is
documented rather than executed.

\vspace{10pt}
\noindent\textbf{Keywords:} embedded systems security; MQTT; configuration defects;
specification mutation; response-time analysis; formal verification; ESP32
